\section*{Bab \ref{ch:g12:stretch-reflex} --- Refleks Regang dan Pesan Saraf}

\begin{solution}{exo:g12:stretch-reflex:1}
Reseptor (kumparan otot), \omterm{def:g12:stretch-reflex:neuron}{neuron} sensorik, pusat pemaduan (materi
kelabu sumsum tulang belakang, satu \omterm{def:g12:stretch-reflex:synapse}{sinapsis}), \omterm{def:g12:stretch-reflex:neuron}{neuron} motorik, dan
efektornya (ototnya).
\end{solution}

\begin{solution}{exo:g12:stretch-reflex:2}
Sebuah badan \omterm{def:g10:cells-common-unit:cell}{sel} dengan intinya, dendrit yang menerima isyarat, dan
satu akson (sering bermielin) yang membawa pesan \omterm{def:g12:stretch-reflex:neuron}{neuron} itu ke
ujungnya. Pesannya berjalan dari dendrit dan badan \omterm{def:g10:cells-common-unit:cell}{sel} sepanjang
aksonnya menuju ujungnya.
\end{solution}

\begin{solution}{exo:g12:stretch-reflex:3}
Pembalikan singkat tegangan membrannya, dari \qty{-70}{mV} menjadi
sekitar \qty{+30}{mV} lalu kembali dalam semilidetik, yang berjalan
sepanjang aksonnya tanpa berubah. Semua atau tak sama sekali: di bawah
ambangnya tak ada apa-apa, di atasnya selalu isyarat penuh yang sama.
Kuatnya tersandi pada frekuensi \omterm{prop:g12:stretch-reflex:actionpotential}{potensial aksinya}.
\end{solution}

\begin{solution}{exo:g12:stretch-reflex:4}
\omterm{prop:g12:stretch-reflex:actionpotential}{Potensial aksi} mencapai ujungnya; vesikelnya melepaskan
\omterm{def:g12:stretch-reflex:synapse}{neurotransmiter} ke dalam celahnya; transmiter itu mengikat reseptor
pada \omterm{def:g10:cells-common-unit:cell}{sel} penerimanya, lalu mengubah tegangannya (menuju atau menjauhi
ambangnya); ia kemudian dihancurkan atau ditarik kembali.
\end{solution}

\begin{solution}{exo:g12:stretch-reflex:5}
Asetilkolina. Kurare menempati reseptornya pada ototnya tanpa
mengaktifkannya; toksin botulinum mencegah pelepasannya dari
vesikelnya. Keduanya melumpuhkan.
\end{solution}

\begin{solution}{exo:g12:stretch-reflex:6}
6, 12 dan 24 paku dalam \qty{100}{ms}: 60, 120 dan 240 tiap detik.
Ukuran dan bentuk pakunya tidak berubah.
\end{solution}

\begin{solution}{exo:g12:stretch-reflex:7}
Waktu penghantarannya $30 - 12 = \qty{18}{ms}$ untuk \qty{1.5}{m}:
sekitar \qty{83}{m/s}.
\end{solution}

\begin{solution}{exo:g12:stretch-reflex:8}
Akar dorsal membawa serat sensoriknya: bila dipotong, pesan regangannya
tak pernah mencapai sumsumnya, tetapi \omterm{def:g12:stretch-reflex:neuron}{neuron} motorik dan seratnya utuh
dan masih dapat dirangsang untuk mengontraksikan ototnya. Akar ventral
membawa serat motoriknya: bila dipotong, sama sekali tak ada yang dapat
mencapai ototnya dari sumsumnya.
\end{solution}

\begin{solution}{exo:g12:stretch-reflex:9}
Bila antagonisnya ikut berkontraksi, ia akan melawan geraknya lalu
meregangkan dirinya sendiri, sehingga memicu \omterm{def:g12:stretch-reflex:reflex}{refleksnya} sendiri ---
sebuah tarik tambang. Interneuron sumsumnyalah yang menghambat, lewat
sebuah \omterm{def:g12:stretch-reflex:synapse}{sinapsis} penghambat pada \omterm{def:g12:stretch-reflex:neuron}{neuron} motorik antagonisnya.
\end{solution}

\begin{solution}{exo:g12:stretch-reflex:10}
1 satuan: tak ada apa-apa (di bawah ambangnya). 4 dan 8 satuan:
\omterm{prop:g12:stretch-reflex:actionpotential}{potensial aksi} sebesar \qty{100}{mV} yang sama, tetapi pada frekuensi
yang lebih tinggi, dan lebih tinggi untuk 8 daripada untuk 4.
\end{solution}

\begin{solution}{exo:g12:stretch-reflex:11}
Hanya ujungnya yang memuat vesikel transmiter dan hanya membran
penerimanya yang memuat reseptor, sehingga zat kimianya hanya dapat
bekerja satu arah. Melepaskan transmiternya, difusinya menyeberangi
celahnya dan tanggapan reseptornya memakan sekitar semilidetik.
\end{solution}

\begin{solution}{exo:g12:stretch-reflex:12}
Asetilkolina menumpuk di sambungannya lalu terus memicu seratnya:
kedutan, lalu kontraksi yang menetap dan kelelahan ototnya, termasuk
otot pernapasannya --- kematian oleh sesak napas. Sebuah penghambat
reseptor (obat semacam atropina) mengurangi rangsangan yang
berlebihan itu; terlalu sedikit membiarkan keracunannya, terlalu banyak
melumpuhkan seperti kurare, jadi dosisnya harus pas.
\end{solution}

\begin{solution}{exo:g12:stretch-reflex:13}
Tiap masukan sensorik merangsang \omterm{def:g12:stretch-reflex:neuron}{neuron} motoriknya dan, dalam keadaan
normal, menghambat \omterm{def:g12:stretch-reflex:neuron}{neuron} antagonisnya lewat interneuron. Dengan
penghambatan yang tersumbat, sebuah sentuhan merangsang kedua sisi
tiap \omterm{def:g10:muscles-and-joints:joint}{sendi} dan rangsangannya menyebar tanpa lawan di dalam sumsumnya:
semua otot berkontraksi sekaligus, dan kejangnya meregangkan otot lain
yang \omterm{def:g12:stretch-reflex:reflex}{refleksnya} menambahi kejang itu.
\end{solution}

\begin{solution}{exo:g12:stretch-reflex:14}
Waktu penghantarannya $60 - 12 = \qty{48}{ms}$ untuk \qty{1.5}{m}:
sekitar \qty{31}{m/s}, sepertiga taraf normal. Mielin membiarkan
\omterm{prop:g12:stretch-reflex:actionpotential}{potensial aksinya} melompat di antara sela pada selubungnya alih-alih
berjalan malar; tanpanya seratnya menghantarkan seperti serat tipis
tanpa mielin.
\end{solution}

\begin{solution}{exo:g12:stretch-reflex:15}
Sumsumnya memang membawa pesan antara otak dan tubuh, tetapi ia juga
mengolahnya: materi kelabunya memuat \omterm{def:g12:stretch-reflex:synapse}{sinapsis} \omterm{def:g12:stretch-reflex:reflex}{lengkung refleksnya},
tempat sebuah regangan diubah menjadi sebuah kontraksi dan, lewat
interneuron, menjadi penghambatan antagonisnya, tanpa satu pesan pun
mencapai otak. Hewan yang sumsumnya diputus dari otaknya masih menarik
kakinya dari cubitan dan masih menunjukkan sentakan lutut: sumsumnya
sebuah pusat, bukan sekadar sebuah kabel.
\end{solution}

\begin{solution}{pb:g12:stretch-reflex:1}
\textbf{1.} $0.75 + 0.75 = \qty{1.5}{m}$.

\textbf{2.} $30 - 4 - 4 - 1 = \qty{21}{ms}$ penghantaran: $1.5/0.021
\approx \qty{70}{m/s}$.

\textbf{3.} Sekitar \qty{10}{cm} lebih tiap arah, \qty{0.2}{m}
seluruhnya: $0.2/70 \approx \qty{3}{ms}$ lebih, jadi sekitar \qty{33}{ms}.

\textbf{4.} Isyarat listriknya adalah pengaktifan seratnya;
kontraksinya sendiri --- meluncurnya filamen dan tarikan pada
\omterm{def:g10:muscles-and-joints:muscle}{tendonnya} --- memakan puluhan milidetik untuk membangun gaya yang cukup
menggerakkan tungkainya.

\textbf{5.} Pesan dari sumsumnya ke otak dan perintah yang kembali
menempuh \qty{1.2}{m} pada sekitar \qty{70}{m/s} --- \qty{17}{ms} ---
ditambah beberapa \omterm{def:g12:stretch-reflex:synapse}{sinapsis} dan pengolahan otaknya sendiri:
sekurang-kurangnya \qty{50}{ms}, dan sebuah keputusan jauh lebih lama.
Tendangannya sudah selesai sebelum isyarat otak mana pun dapat mencapai
\omterm{def:g12:stretch-reflex:neuron}{neuron} motoriknya.

\textbf{6.} Pada 200 tiap detik, 2 \omterm{prop:g12:stretch-reflex:actionpotential}{potensial aksi} dalam \qty{10}{ms}
(0.2 saat istirahat). Lajunya berubah; ukuran dan bentuk tiap pakunya
tidak.

\textbf{7.} Sebesar 10 kali lipat.

\textbf{8.} Pada 20 tiap detik, dosis transmiternya berjarak
\qty{50}{ms} dan masing-masing memudar sebelum yang berikutnya datang;
\omterm{def:g12:stretch-reflex:neuron}{neuron} motoriknya tak pernah mencapai ambangnya. Pada 200 tiap detik
dosisnya tiba berjarak \qty{5}{ms} lalu menjumlah: ambangnya terlewati
dan \omterm{def:g12:stretch-reflex:neuron}{neuron} motoriknya pun melepaskan isyarat.

\textbf{9.} Lebih sedikit \omterm{def:g12:stretch-reflex:neuron}{neuron} motorik yang mencapai ambangnya dan
masing-masing melepaskan lebih sedikit \omterm{prop:g12:stretch-reflex:actionpotential}{potensial aksi}: lebih sedikit
unit motorik yang berkontraksi, pada laju yang lebih rendah, dan
tendangannya pun lebih lemah.

\textbf{10.} \omterm{prop:g12:stretch-reflex:actionpotential}{Potensial aksi} bersifat semua atau tak sama sekali:
besarnya ditetapkan oleh \omterm{def:g12:stretch-reflex:neuron}{neuronnya}, sehingga besar itu tak membawa
keterangan tentang rangsangannya. Hanya lajunya yang dapat berubah.

\textbf{11.} Asetilkolina; setelah mengikat reseptornya, ia dihancurkan
dalam beberapa milidetik oleh sebuah \omterm{def:g11:enzymes-and-phenotype:enzyme}{enzim} di dalam celahnya.

\textbf{12.} Reseptor \omterm{def:g10:muscles-and-joints:muscle}{serat ototnya}: asetilkolinanya dilepaskan tetapi
tak dapat bekerja.

\textbf{13.} Pelepasan asetilkolina dari vesikelnya. Di laboratorium,
asetilkolina yang diterapkan langsung pada ototnya akan membuatnya
berkontraksi di bawah toksin botulinum (reseptornya bebas) tetapi tidak
di bawah kurare (reseptornya terhambat).

\textbf{14.} \omterm{prop:g12:stretch-reflex:actionpotential}{Potensial aksi} serat sensoriknya dan, karena tak ada yang
mencapai sumsumnya, juga \omterm{prop:g12:stretch-reflex:actionpotential}{potensial aksi} \omterm{def:g12:stretch-reflex:neuron}{neuron} motorik serta
elektromiogramnya; ototnya masih berkontraksi bila sarafnya dirangsang
di bawah sumbatannya.

\textbf{15.} Kurare menyumbat penerimaannya, toksin botulinum
menyumbat pelepasannya, anestetiknya menyumbat penghantarannya: tiga
langkah, satu otot yang senyap.

\textbf{16.} Serat sensoriknya juga merangsang sebuah interneuron yang
menghambat \omterm{def:g12:stretch-reflex:neuron}{neuron} motorik otot paha belakangnya; kegiatan istirahatnya
ditekan, dan antagonisnya mengendur di bawah tonus istirahatnya.

\textbf{17.} \omterm{def:g12:stretch-reflex:reflex}{Refleksnya} hanya memerlukan sumsumnya, yang masih utuh
pada aras itu. Berlebihannya menunjukkan bahwa otak biasanya mengirim
isyarat yang meredam \omterm{def:g12:stretch-reflex:reflex}{refleks} itu; tanpanya lengkungnya berjalan pada
kepekaan penuh.

\textbf{18.} Kumparannya, saraf sensorik atau akar dorsalnya, materi
kelabu sumsumnya pada aras itu, saraf motorik atau akar ventralnya,
atau ototnya sendiri. Merangsang saraf motoriknya langsung lalu merekam
ototnya memberi tahu apakah sisi motoriknya bekerja; merekam saraf
sensoriknya selama sebuah regangan memberi tahu apakah sisi sensoriknya
bekerja.

\textbf{19.} \omterm{def:g12:stretch-reflex:reflex}{Refleks} itu menguji, dalam satu ketukan, seluruh
lengkungnya --- saraf sensorik, sumsum, saraf motorik, otot --- dan
peredaman oleh otaknya: \omterm{def:g12:stretch-reflex:reflex}{refleks} yang hilang, lemah atau berlebihan
menunjuk sebuah aras sistem sarafnya sebelum tanda lain mana pun.

\textbf{20.} Sekitar \qty{70}{m/s}; frekuensi \omterm{prop:g12:stretch-reflex:actionpotential}{potensial aksi}
kumparannya; dan pengikatan asetilkolina pada reseptor ototnya.
\end{solution}
